\documentclass[11pt]{amsart}
\usepackage{amsmath,amssymb,amsthm}

\newtheorem{theorem}{Theorem}
\newtheorem{proposition}[theorem]{Proposition}
\newtheorem{corollary}[theorem]{Corollary}

\title{A characterization of 6 via arithmetic functions}
\author{}
\date{\today}

\begin{document}
\maketitle

\begin{abstract}
We prove that $\sigma(n)\varphi(n) = n\tau(n)$ if and only if $n \in \{1,6\}$,
providing a new characterization of the smallest perfect number using only the
four basic arithmetic functions $\sigma$, $\varphi$, $\tau$, and the identity.
The dual equation $\sigma(n)\tau(n) = n\varphi(n)$ characterizes $n=28$.
We establish a spectral gap: the ratio $R(n) = \sigma(n)\varphi(n)/(n\tau(n))$
lies in $\{3/4\} \cup \{1\} \cup [7/6, \infty)$, with both intervals $(3/4,1)$
and $(1,7/6)$ empty.
\end{abstract}

\section{Introduction}

Let $\sigma(n)$, $\tau(n)$, and $\varphi(n)$ denote the sum of divisors,
number of divisors, and Euler's totient function, respectively. Define
\[
  R(n) = \frac{\sigma(n)\varphi(n)}{n\,\tau(n)}.
\]

\begin{theorem}\label{thm:main}
$R(n) = 1$ if and only if $n \in \{1, 6\}$.
\end{theorem}

\begin{theorem}\label{thm:gap}
For all $n \geq 2$,
\[
  R(n) \in \left\{\tfrac{3}{4}\right\} \cup \{1\} \cup \left[\tfrac{7}{6}, \infty\right).
\]
In particular, both $(3/4, 1)$ and $(1, 7/6)$ are empty.
\end{theorem}

\begin{theorem}\label{thm:dual}
$\sigma(n)\tau(n) = n\varphi(n)$ if and only if $n \in \{1, 28\}$.
\end{theorem}

\section{Proof of Theorem~\ref{thm:main}}

Since $\sigma$, $\varphi$, $\tau$ are multiplicative and $n \mapsto n$ is
completely multiplicative, $R$ is multiplicative on squarefree integers and
factors as
\[
  R(n) = \prod_{p^a \| n} f(p,a), \quad
  f(p,a) = \frac{p^{a+1}-1}{p(a+1)}.
\]

\textbf{Step 1.} $f(2,1) = 3/4$ is the unique value below 1.

For $a=1$: $f(p,1) = (p^2-1)/(2p) \geq 4/3$ when $p \geq 3$.
For $a \geq 2$: $f(p,a) \geq f(2,2) = 7/6 > 1$.

\textbf{Step 2.} $R(n)=1$ requires $2 \| n$.

If $2 \nmid n$, all factors exceed 1. If $4 \mid n$, the 2-part contributes
$f(2,a) \geq 7/6 > 1$.

\textbf{Step 3.} The remaining product equals $4/3$.

We need $f(q,1) = 4/3$, giving $(q^2-1)/(2q) = 4/3$, so $3q^2 - 8q - 3 = 0$,
with unique prime solution $q = 3$.

\textbf{Step 4.} No further factors allowed.

Any additional $f(r,b) \geq 4/3$ pushes the product above $4/3$.
Hence $n = 2 \times 3 = 6$.  \qed

% TODO: Complete sections 3-8

\end{document}
